% ID: SEL-D0077-OB-E03
% Tanda S077.
\textbf{Parte I.} Sean $O$ disponer de ordenador y $T$ poder ver la TDT.

\textbf{a)} Se conoce $P(T\mid O)=0{,}80$, que difiere de $P(T)=0{,}60$. Como $P(O)>0$, esta desigualdad prueba que los sucesos son dependientes. Equivalentemente, $P(O\cap T)=0{,}70\cdot0{,}80=0{,}56$ difiere de $P(O)P(T)=0{,}42$.

\textbf{b)} Usando inclusión y exclusión, $P(O\cup T)=0{,}70+0{,}60-0{,}56=0{,}74$. El complementario es no tener ninguno de los dos equipamientos. Su probabilidad es $1-0{,}74=0{,}26$, por lo que el porcentaje pedido es $\boxed{26\%}$.

\textbf{Parte II.} La proporción muestral es $\hat p=60/250=0{,}24$. Hay 60 portadoras y 190 no portadoras, por lo que usamos la aproximación normal, suponiendo que la muestra es aleatoria y representativa.

Al $97\%$ queda $0{,}015$ en cada cola y $z_{0{,}985}\approx2{,}1701$. Aplicamos el intervalo para una proporción:
\[
 0{,}24\pm2{,}1701\sqrt{\frac{0{,}24\cdot0{,}76}{250}}
 \approx\boxed{[0{,}1814;\ 0{,}2986]}.
\]
